%%%% CAPÍTULO 1 - INTRODUÇÃO
%%
%% Deve apresentar uma visão global da pesquisa, incluindo: breve histórico, importância e justificativa da escolha do tema,
%% delimitações do assunto, formulação de hipóteses e objetivos da pesquisa e estrutura do trabalho.

%% Título e rótulo de capítulo (rótulos não devem conter caracteres especiais, acentuados ou cedilha)
\chapter{Introdução}\label{cap:introducao}

\textcolor{red}{João Victor: 1) traduzir todas figuras, procurar manter os nomes dos elementos PON em inglês. 2) Escrever as seções de big data, docker e kubernetes. Clayton: Adaptar a introdução para o contexto de: Investigar a possibilidade de distribuir uma aplicação do paradigma orientado a notificações.}

O aumento contínuo da produção e do consumo de dados em formato de vídeo, aliado à necessidade de responder a consultas sobre eventos em tempo real, tem impulsionado o desenvolvimento de novas abordagens para processamento de dados em fluxo \cite{kossoski2025slr}. Nesse cenário, surgem desafios importantes relacionados ao custo computacional, à latência de resposta e à organização das etapas de ingestão, processamento e recuperação de informação \cite{kossoski2024mdpinopquery}. Em especial, aplicações que lidam com eventos em vídeo tendem a exigir estruturas complexas para identificar objetos, relações espaciais e padrões temporais, o que torna o problema relevante tanto do ponto de vista acadêmico quanto prático \cite{kossoski2024mdpinopquery}.

A literatura recente aponta que, apesar dos avanços em computação, visão computacional e arquiteturas de processamento, ainda existem limitações quando o objetivo é consultar eventos em fluxos contínuos com baixo atraso e uso controlado de recursos \cite{kossoski2024mdpinopquery}. Parte dessas limitações está associada à forma como os sistemas são estruturados, muitas vezes com forte acoplamento entre componentes e com avaliações repetitivas que elevam o custo da execução \cite{kossoski2024mdpinopquery}. 
Ademais, o elevado custo dos componentes de hardware, especialmente das memórias de armazenamento primário (RAM) e dos dispositivos de armazenamento secundário, como discos rígidos (HDs) e unidades de estado sólido (SSDs), tem dificultado a ampliação da infraestrutura computacional tanto de empresas quanto de usuários domésticos \cite{semensato2026ram}. 
Nesse contexto, torna-se ainda mais relevante o desenvolvimento de soluções computacionais eficientes. Corroborando essa perspectiva, John Carmack, reconhecido como um dos mais influentes desenvolvedores de jogos eletrônicos da história, argumenta que a adequada otimização do software pode prolongar a vida útil e ampliar a capacidade de aproveitamento do hardware disponível, reduzindo a necessidade de constantes atualizações de equipamentos \cite{plaza2026doom}. Embora a afirmação tenha sido formulada no contexto da indústria de jogos eletrônicos, seus princípios são amplamente aplicáveis ao desenvolvimento de software como um todo, uma vez que a otimização de sistemas constitui um fator determinante para o melhor aproveitamento dos recursos computacionais disponíveis.


Diante desse cenário, arquiteturas distribuídas ou baseadas em contêineres como Docker e Kubernetes têm se consolidado como alternativas promissoras para a organização de aplicações computacionais complexas \cite{docker2026documentation, kubernetes2026documentation}. Essas abordagens favorecem a separação de responsabilidades entre os diferentes componentes do sistema, promovem maior modularidade e facilitam o gerenciamento da infraestrutura. Além disso, oferecem uma base adequada para a execução de cargas de trabalho intensivas, contribuindo para a escalabilidade, o desempenho e a utilização mais eficiente dos recursos computacionais disponíveis.


É nesse contexto que se insere o Paradigma Orientado a Notificações (PON). Segundo \citet{Neves2021mastersframework4}, o PON propõe uma organização de software baseada em entidades colaborativas e fracamente acopladas, que interagem por meio de notificações direcionadas entre atributos, premissas, condições, regras, ações e métodos. Essa abordagem promove uma estrutura mais modular e reativa, na qual os componentes do sistema respondem a eventos de forma coordenada e eficiente. Além disso, busca reduzir redundâncias estruturais e temporais, favorecendo o reaproveitamento de recursos computacionais e a simplificação da lógica de execução. Como consequência, o paradigma apresenta potencial para facilitar a distribuição do processamento e melhorar a escalabilidade de aplicações em cenários que demandam elevado desempenho computacional \cite{ronszcka2019phdcompiladores}. 

Por sua vez, o trabalho de \citet{kossoski2024mdpinopquery} adaptou os conceitos e ferramenta de desenvolvimento do PON ao contexto de consultas em fluxos contínuos de vídeo, com o objetivo de reduzir a latência de processamento e minimizar a dependência de operações intensivas em bancos de dados. Além de apresentar resultados promissores para aplicações de \textit{video querying}, os autores destacam a possibilidade de evolução da proposta para ambientes distribuídos, especialmente em cenários caracterizados por grandes volumes de dados e elevadas demandas computacionais \cite{kossoski2024mdpinopquery}.

O presente Trabalho de Conclusão de Curso pesquisa retoma essa perspectiva, porém com um escopo mais específico. Em vez de propor novos mecanismos de consulta ou aprofundar aspectos relacionados à visão computacional ou banco de dados, o trabalho concentra-se na investigação da viabilidade de execução de uma aplicação baseada no PON em uma infraestrutura distribuída. Para esse propósito, este trabalho propõe o desenvolvimento de uma prova de conceito implantada em um cluster Kubernetes, utilizando contêineres Docker para o empacotamento dos serviços e Kubernetes como plataforma de orquestração \cite{kubernetes2026documentation}\cite{docker2026documentation}.

Dessa forma, o foco do trabalho desloca-se para a análise da organização arquitetural da aplicação, contemplando a distribuição de seus componentes, sua execução em ambiente conteinerizado e distribuído, bem como a observação de características que possam subsidiar futuras avaliações relacionadas ao desempenho, à escalabilidade e à eficiência na utilização dos recursos computacionais.


\section{Problema}

O problema que fundamenta esta pesquisa não consiste simplesmente em executar uma aplicação em um ambiente Kubernetes, mas em investigar a viabilidade de organizar, conteinerizar e distribuir uma aplicação de prova de conceito baseada no PON em um cluster Kubernetes. A motivação para essa investigação decorre do fato de que o PON foi concebido para operar com entidades colaborativas e potencialmente distribuídas, porém, até o momento, não foram identificados trabalhos que explorem sua implantação em ambientes orquestrados por Kubernetes.

Embora o PON tenha sido proposto como uma abordagem voltada à redução de redundâncias estruturais e temporais, bem como à promoção de uma execução reativa e desacoplada, sua aplicação prática ainda demanda validação em diferentes contextos arquiteturais \cite{kossoski2024mdpinopquery}. Nesse sentido, esta pesquisa parte da hipótese de que a distribuição dos componentes da aplicação em um cluster Kubernetes pode proporcionar melhores condições para observar aspectos relacionados à separação de responsabilidades, à comunicação entre serviços, ao isolamento proporcionado pelos contêineres e à possibilidade de escalabilidade seletiva dos componentes do sistema.

Entretanto, não se assume que a simples adoção do Kubernetes seja suficiente para solucionar eventuais limitações de desempenho ou reduzir, por si só, o custo computacional do processamento. Em vez disso, o Kubernetes é compreendido como uma plataforma que possibilita a organização da aplicação em uma infraestrutura distribuída, fornecendo mecanismos que podem facilitar futuras análises de desempenho, escalabilidade, disponibilidade e utilização de recursos computacionais \cite{kubernetes2026documentation}.

Diante desse contexto, a questão que orienta esta pesquisa pode ser formulada da seguinte maneira: é viável organizar, conteinerizar e distribuir uma aplicação de prova de conceito baseada no PON em um cluster Kubernetes, de modo a possibilitar sua execução em ambiente distribuído e estabelecer uma infraestrutura adequada para futuras avaliações de desempenho e escalabilidade?


\section{Objetivos geral e específicos}

\subsection{Objetivo Geral}

Investigar a viabilidade da distribuição de uma aplicação de prova de conceito baseada no Paradigma Orientado a Notificações (PON) em um cluster Kubernetes, por meio da conteinerização e da orquestração de seus componentes, visando analisar sua organização arquitetural, sua execução em ambiente distribuído e seu potencial para futuras avaliações de desempenho e escalabilidade.

\subsection{Objetivos Específicos}

\begin{itemize}
\item Estudar os fundamentos teóricos do Paradigma Orientado a Notificações e as características da aplicação utilizada como prova de conceito.
\item Analisar a arquitetura da aplicação, identificando seus componentes, dependências e mecanismos de comunicação.
\item Conteinerizar os componentes da aplicação utilizando a plataforma Docker.
\item Projetar e implementar a infraestrutura necessária para a execução da aplicação em um cluster Kubernetes.
\item Configurar os recursos de orquestração necessários, incluindo \textit{Deployments}, \textit{Services}, \textit{ConfigMaps} e demais artefatos pertinentes.
\item Implantar e executar a aplicação em ambiente distribuído.
\item Avaliar aspectos relacionados à distribuição dos componentes, à comunicação entre serviços, ao isolamento proporcionado pelos contêineres e à organização arquitetural da solução.
\item Identificar limitações, desafios e oportunidades observados durante o processo de implantação e execução.
\item Documentar os resultados obtidos e discutir a viabilidade da abordagem proposta como base para futuras investigações sobre desempenho e escalabilidade.
\end{itemize}


\section{Justificativa}

A relevância desta pesquisa pode ser compreendida, em primeiro lugar, pela continuidade de uma linha de investigação já estabelecida em torno do Paradigma Orientado a Notificações (PON). Os trabalhos desenvolvidos até o momento forneceram uma base conceitual e experimental consistente para o paradigma, demonstrando seu potencial em diferentes contextos de aplicação. Entretanto, ainda existe uma lacuna relacionada à sua utilização em ambientes distribuídos e conteinerizados, aspecto que permanece pouco explorado na literatura. Nesse sentido, ao dar continuidade a essa linha de pesquisa, o presente trabalho contribui para ampliar o entendimento sobre a aplicabilidade do PON em cenários alinhados às arquiteturas de software contemporâneas.

Sob a perspectiva prática, a transformação da aplicação de prova de conceito baseada em PON em uma solução conteinerizada e distribuída amplia significativamente seu potencial como plataforma experimental. A utilização do Docker favorece a padronização e a portabilidade do ambiente de execução, enquanto o Kubernetes oferece mecanismos para gerenciamento, replicação, escalonamento e orquestração dos serviços que compõem a aplicação \cite{docker2026documentation}\cite{kubernetes2026documentation}. Dessa forma, a infraestrutura desenvolvida pode servir como base para futuras investigações envolvendo desempenho, escalabilidade, balanceamento de carga, tolerância a falhas e comparações entre diferentes estratégias arquiteturais.

Adicionalmente, o Kubernetes tem se consolidado como uma das principais plataformas para implantação e gerenciamento de aplicações distribuídas baseadas em contêineres, sendo amplamente adotado em ambientes acadêmicos e corporativos que demandam flexibilidade, disponibilidade e escalabilidade operacional. Nesse contexto, a análise de sua aplicação em uma solução fundamentada no PON promove a integração de temas centrais da Engenharia da Computação, tais como arquitetura de software, sistemas distribuídos, computação em nuvem, conteinerização e orquestração de serviços. Consequentemente, além de sua contribuição técnica, esta pesquisa possui relevância formativa, ao articular fundamentos teóricos e tecnológicos em uma aplicação prática voltada à investigação de problemas reais de interesse científico.


\section{Estrutura do trabalho}

Este trabalho está organizado em capítulos que acompanham a evolução da investigação proposta. O Capítulo 2 apresenta a fundamentação teórica, abordando o Paradigma Orientado a Notificações, o NOP Query, além de conceitos relacionados a aplicações distribuídas, Docker e Kubernetes. Essa base é necessária para situar o leitor e justificar as escolhas feitas ao longo do desenvolvimento.

O Capítulo 3 descreve a metodologia adotada, indicando os procedimentos empregados para preparar a aplicação prova de conceito, conteinerizar seus componentes e implantar a solução no cluster Kubernetes. Nessa etapa, são apresentados os critérios técnicos e as decisões que orientam a estruturação do ambiente distribuído.

No Capítulo 4, são descritas as atividades de desenvolvimento, os experimentos realizados e as observações obtidas durante a execução da aplicação. Esse capítulo concentra a análise da viabilidade da abordagem, considerando funcionamento, organização dos serviços, comunicação entre componentes e limitações encontradas.

Por fim, o Capítulo 5 apresenta as considerações finais, retomando os objetivos definidos no início do trabalho, sintetizando os resultados obtidos e destacando limitações da proposta. Também são indicadas possibilidades de continuidade da pesquisa, especialmente em relação a testes de desempenho, escalabilidade e aprimoramentos arquiteturais.

\section{Cronograma}

O desenvolvimento deste trabalho foi organizado em etapas, de modo a permitir a progressão gradual das atividades de pesquisa, implementação e escrita. Essa organização é importante porque a investigação envolve desde o estudo conceitual até a preparação do ambiente distribuído e a análise dos resultados obtidos. A tabela a seguir apresenta um cronograma genérico, com períodos indicados em meses, que pode ser adaptado conforme a evolução real do TCC.

\begin{table}[htbp]
\centering
\caption{Cronograma genérico de desenvolvimento do TCC}
\label{tab:cronograma}

\small
\renewcommand{\arraystretch}{1.3}

\begin{tabular}{p{7.2cm}cccccccc}
\hline
\textbf{Etapa} &
\textbf{M1} &
\textbf{M2} &
\textbf{M3} &
\textbf{M4} &
\textbf{M5} &
\textbf{M6} &
\textbf{M7} &
\textbf{M8} \\
\hline

Revisão bibliográfica sobre NOP, NOP Query, Docker e Kubernetes
& X & X & & & & & & \\

Estudo da aplicação prova de conceito
& X & X & X & & & & & \\

Preparação do ambiente de desenvolvimento
& & X & X & & & & & \\

Conteinerização da aplicação
& & & X & X & & & & \\

Configuração do cluster Kubernetes
& & & & X & X & & & \\

Implantação e testes iniciais
& & & & & X & X & & \\

Coleta de observações e resultados
& & & & & & X & X & \\

Escrita e revisão do TCC
& X & X & X & X & X & X & X & X \\

\hline
\end{tabular}
\end{table}
% Segundo \citeonline{Coulouris2013}.

% Segundo \citeonline[p. 40]{Coulouris2013}.

% Citação no final do Parágrafo~\cite{Coulouris2013}. 

% Citação no final do Parágrafo com número de página~\cite[p. 40]{Coulouris2013}.

% %(Modelo de referência: pessoa jurídica)
% Citação no final do Parágrafo~\cite{NBR6023:2018}

% %(Modelo de referência: pessoa jurídica)
% Citação no final do Parágrafo~\cite{NBR6027:2012}

% %(Modelo de referência: pessoa jurídica)
% Citação no final do Parágrafo~\cite{NBR6028:2021}

% Segundo a \citeonline{NBR14724:2011}.

% Citação no final do Parágrafo~\cite{NBR10520:2002}

% Citação no final do Parágrafo~\cite{NBR14724:2011}.

% % (Modelo de referência de trabalho acadêmico).
% Citação no final do Parágrafo~\cite{Andrade2005}

% % (Modelo de referência: capítulo de livro).
% Citação no final do Parágrafo~\cite{Borges2014}

% % (Modelo de referência: leis, decretos, portarias, etc.)
% Citação no final do Parágrafo~\cite{BRASIL:1998}

% % (Modelo de referência: livro com subtítulo). Nome com sufixo "Von" - Configuração no bib
% Citação no final do Parágrafo~\cite[p. 66]{KROGH:2001}

% Citação no final do Parágrafo~\cite{Faina2001}

% % (Modelo de referência: livro com subtítulo).
% Citação no final do Parágrafo~\cite{Davenport2012}

% % (Modelo de referência: artigo de periódico).
% Citação no final do Parágrafo~\cite{Monteiro2009}

% %(Modelo de referência: artigo de periódico). Nome familiar "Junior"
% Citação no final do Parágrafo~\cite{Sanches2024}

% % (Modelo de referência: trabalho publicado em evento).
% Citação no final do Parágrafo~\cite{Renaux2001}

\section{Contextualização}

Aplicações que processam dados em fluxo têm se tornado cada vez mais relevantes em domínios que dependem de resposta rápida, como monitoramento, análise de eventos e consulta a conteúdo multimídia. No caso de vídeos, o problema se torna ainda mais complexo porque os dados são volumosos, dinâmicos e compostos por múltiplos elementos semiestruturados, como objetos, atributos e relações espaciais e temporais \cite{kossoski2024mdpinopquery}. Nesses cenários, a consulta não se limita à recuperação de registros armazenados, mas envolve a interpretação de eventos à medida que os dados chegam ao sistema, o que amplia a exigência sobre processamento, comunicação entre módulos e consumo de recursos.

Kossoski, Simão e Lopes destacam que o crescimento dos dados de vídeo torna a busca por eventos em vídeos cada vez mais relevante, ao mesmo tempo em que reforça a presença de desafios relacionados ao custo computacional e à latência \cite{kossoski2024mdpinopquery}. Esses desafios aparecem tanto em abordagens que utilizam bancos de dados e estruturas tradicionais quanto em soluções voltadas ao processamento de fluxo contínuo. Quando consultas mais complexas são executadas, especialmente aquelas que envolvem comparação entre objetos, relações espaciais e janelas de tempo, o sistema pode exigir muitas avaliações internas, o que impacta diretamente o tempo de resposta \cite{kossoski2024mdpinopquery}.

Dentro desse contexto, o NOP Query é apresentado como um método orientado a notificações para processamento de buscas em vídeos em tempo real, apoiado nos princípios do Paradigma Orientado a Notificações \cite{kossoski2024mdpinopquery}. A principal característica dessa abordagem é a organização do processamento em pequenas entidades colaborativas, que reagem a mudanças por meio de notificações bem definidas, evitando verificações desnecessárias e reduzindo acoplamentos \cite{kossoski2024mdpinopquery}. Em vez de depender de uma estrutura centrada em buscas intensivas ou em avaliações repetitivas, o método procura reagir ao surgimento de eventos relevantes, o que o torna conceitualmente atrativo para aplicações em fluxo.

A tese de referência também sinaliza interesse em investigar arquitetura distribuída para uma aplicação de big data no domínio de video querying, utilizando ferramentas do ecossistema NOP \cite{kossoski2024mdpinopquery}. Esse ponto é importante porque a distribuição não é apenas uma decisão de infraestrutura; ela afeta a forma como os componentes se comunicam, como as tarefas são divididas e como a aplicação responde sob carga. Assim, avaliar a possibilidade de execução distribuída em cluster torna-se um passo coerente para compreender melhor o comportamento do método em um ambiente mais próximo de aplicações modernas.

Nesse cenário, Docker e Kubernetes assumem papel central na proposta deste TCC. Docker permite empacotar a aplicação e suas dependências em contêineres, favorecendo reprodutibilidade, isolamento e portabilidade entre ambientes \cite{docker2024}. Kubernetes, por sua vez, atua como plataforma de orquestração de contêineres, oferecendo mecanismos para implantação, escalonamento, gerenciamento de serviços e reorganização da aplicação em um cluster \cite{kubernetes2024}. Dessa forma, essas tecnologias viabilizam a montagem de um ambiente distribuído no qual a aplicação prova de conceito pode ser executada de forma organizada, permitindo observar aspectos arquiteturais e operacionais relevantes para a investigação proposta.

